% Appendix: EMF chain in detail. Expands sec:emf; \input by report_2.tex.
\section{EMF chain in detail}\label{app:emf}

This appendix expands \cref{sec:emf}: how the EMF chain turns one \texttt{.exam} file into a Python
script in three model-driven steps, run headless by an Eclipse application; the script then solves the
Problem and validates what it finds. The models parsed from the file, and the Conflict Graph, are checked
against their metamodels' multiplicities and OCL invariants (\cref{app:metamodels}) before the next step
uses them.

\paragraph{Xtext.} The grammar imports the hand-written Problem and Solver Configuration metamodels, so one
\texttt{.exam} file holds a Problem (its batches, its students, and each exam with its optional weight, 1 by
default, and the students who sit it) and, optionally, its Solver Configuration: the Method, the Method's
settings, the seed and the Penalty Coefficient, automatic unless given (\cref{app:instances}).

\paragraph{ATL.} An EMFTVM transformation makes the Problem a Conflict Graph \cite[§5.2, Fig. 3]{Paper}. It
copies the exams with their weights and the number of batches~$M$, the range of the index~$k$ of $x_{ik}$
\cite[§5.3, Eq. (4)]{Paper}, and adds one Conflict for each pair of exams that share a student, the
paper's course conflict matrix~$C$ \cites[§5.2, Algorithm 1, Eq. (3)]{Paper}[cell 7]{Notebook}.

\paragraph{Acceleo.} A template generates \texttt{solve\_<name>.py} from the Conflict Graph and the Solver
Configuration: both as Python data, then the notebook's pyqubo construction
\cites[cells 11--13]{Notebook}{pyqubo} (\texttt{emf/docs/adr/0007}), made the weighted set packing of
\cref{sec:strategy} (\texttt{emf/docs/adr/0009}). Each
Constraint is labelled in domain terms, \texttt{at-most-one-batch-<exam>} or
\texttt{conflict-<exam>-<exam>-batch-<k>}, so a broken one names its exams. pyqubo compiles $H$ into the
QUBO and its Ising form, so the chain models neither.

\paragraph{The runtime.} The generated script holds only what depends on the models; it imports the
hand-written \texttt{emf/python/examschedule.py}, which holds the rest (\texttt{emf/docs/adr/0005}). It sets
$A$ to the largest weight plus one, or rejects an explicit $A$ not above every weight
\cite[§4, Comment on the Scalar Penalty P]{Glover}, and samples by the Method chosen. Simulated annealing
runs dwave-neal on the compiled model, as the notebook does \cite[cell 15]{Notebook}. QAOA \cite{Farhi}
builds the QAOA Circuit from the Ising form, with $p$ alternating cost and mixer layers on the uniform
superposition \cite[Eqs. (2)--(6)]{Farhi}, runs it with Qiskit \cite{Qiskit} on the Aer simulator, and lets a classical
optimiser tune the $2p$ angles to minimise the expected energy (Farhi et al.\ maximise theirs)
\cite[Eqs. (7)--(8)]{Farhi}, starting on the annealing path \cite[§VI]{Farhi}. The runtime then decodes
the lowest-energy sample into a Schedule model \cite[§5.5]{Paper}, runs the Baseline, and writes the run's
results, from which \cref{tab:results-emf} is generated.

Every ATL rule, and every template section and runtime step that implements the formulation or QAOA, cites
its reference in a comment beside it, as
\texttt{[Key locator]}; \cref{tab:traceability} is generated from those comments.
